\section{Сплайн-аппроксимация}
\label{sec:spline-approximation}

\subsection{Определение сплайна}

Пусть задано разбиение
\[
    a=x_0<x_1<\cdots<x_n=b.
\]
Сплайном степени $p$ называется функция, которая на каждом отрезке
$[x_{i-1},x_i]$ является многочленом степени не выше $p$ и обладает заданной
гладкостью в узлах. Наиболее распространены кубические сплайны класса $C^2$.

Главное преимущество сплайна перед глобальным многочленом --- локальность:
изменение одного узла влияет в основном на соседние участки, а степень остаётся
малой независимо от числа данных.

\subsection{Интерполирующий кубический сплайн}

Кубический сплайн $S$ должен удовлетворять:
\begin{enumerate}
    \item $S_i$ --- кубический многочлен на $[x_{i-1},x_i]$;
    \item $S(x_i)=y_i$;
    \item $S$, $S'$ и $S''$ непрерывны во внутренних узлах;
    \item заданы два дополнительных граничных условия.
\end{enumerate}

Обозначим
\[
    h_i=x_i-x_{i-1},
    \qquad
    M_i=S''(x_i).
\]
На отрезке $[x_{i-1},x_i]$ сплайн можно записать как
\[
\begin{aligned}
S_i(x)=&
\frac{M_{i-1}(x_i-x)^3}{6h_i}
+\frac{M_i(x-x_{i-1})^3}{6h_i}\\
&+\left(y_{i-1}-\frac{M_{i-1}h_i^2}{6}\right)
  \frac{x_i-x}{h_i}
+\left(y_i-\frac{M_i h_i^2}{6}\right)
  \frac{x-x_{i-1}}{h_i}.
\end{aligned}
\]
Непрерывность первой производной приводит к трёхдиагональной системе
\[
    h_iM_{i-1}+2(h_i+h_{i+1})M_i+h_{i+1}M_{i+1}
    =
    6\left(
    \frac{y_{i+1}-y_i}{h_{i+1}}
    -\frac{y_i-y_{i-1}}{h_i}
    \right),
\]
$i=1,\ldots,n-1$.

\subsection{Граничные условия}

Для замыкания системы используют два условия:
\begin{itemize}
    \item естественный сплайн:
    \[
        S''(a)=S''(b)=0;
    \]
    \item закреплённый сплайн:
    \[
        S'(a)=f'_a,
        \qquad
        S'(b)=f'_b;
    \]
    \item периодический сплайн: совпадают значения производных на концах;
    \item условие not-a-knot: первые два и последние два участка образуют по
          одному кубическому многочлену.
\end{itemize}

При естественных условиях система для $M_i$ имеет строгую диагональную
преобладаемость и решается методом прогонки за $O(n)$.

\subsection{Вариационное свойство}

Естественный интерполирующий кубический сплайн минимизирует функционал изгиба
\[
    J[v]=\int_a^b\bigl(v''(x)\bigr)^2\,dx
\]
среди функций $v\in H^2(a,b)$, проходящих через все заданные точки.
Это объясняет его геометрическую гладкость: среди интерполянтов он имеет
минимальную ``кривизну'' в интегральном смысле.

\subsection{Оценки погрешности}

Если $f\in C^4[a,b]$, сетка квазиравномерна, а $h=\max h_i$, то для кубического
интерполирующего сплайна типичны оценки
\[
    \|f-S\|_{\infty}=O(h^4),
    \qquad
    \|f'-S'\|_{\infty}=O(h^3),
    \qquad
    \|f''-S''\|_{\infty}=O(h^2).
\]
Точные константы зависят от граничных условий и регулярности сетки.

\subsection{Сглаживающий сплайн}

Для зашумлённых данных точная интерполяция нежелательна. Сглаживающий кубический
сплайн минимизирует
\[
    J[S]=
    \sum_{i=0}^{n}w_i\bigl(y_i-S(x_i)\bigr)^2
    +\lambda\int_a^b\bigl(S''(x)\bigr)^2\,dx,
    \qquad \lambda\geq0.
\]
Первое слагаемое отвечает за близость к данным, второе --- за гладкость.
При $\lambda\to0$ получается интерполирующий сплайн, при
$\lambda\to\infty$ решение приближается к линейной регрессии.

Параметр $\lambda$ выбирают по оценке шума, перекрёстной проверке или обобщённой
кросс-валидации.

\subsection{$B$-сплайны}

$B$-сплайны образуют локальный базис пространства сплайнов. Сплайн записывается
как
\[
    S(x)=\sum_j c_jB_{j,p}(x).
\]
Каждая базисная функция $B_{j,p}$ отлична от нуля лишь на небольшом числе
соседних интервалов. Это даёт разреженные матрицы, локальное управление формой и
численную устойчивость. Базис строится рекурсией Кокса--де Бура.

\subsection{Что важно сказать на экзамене}

Нужно определить кусочно-полиномиальную структуру и гладкость, вывести или хотя
бы записать трёхдиагональную систему для вторых производных кубического сплайна,
перечислить граничные условия и объяснить отличие интерполирующего сплайна от
сглаживающего.
